\begin{center}
{\Large\bfseries Robotic Active Matter: A Living Systematic Review
of Experimental Platforms, Collective Phenomena, and Tests of Theory}
\end{center}

\section*{Proposal to the Editors of Reviews of Modern Physics}

\noindent\textbf{Proposed article type:}
Review article, maintained as a living systematic review.

\noindent\textbf{Authors:}
Andreas Maier, Siming Bayer, M. Delens, Isa Mammadli,
M. M\'elard, Siddhant Mohapatra, Maxime Noirhomme,
Filip Novkoski, Jayant Pande, Prajol Shrestha,
Ana-Sun\v{c}ana Smith, N. Vanesse, F. W\'ery,
and Nicolas Vandewalle.

\noindent\textbf{Affiliations:}

\noindent Andreas Maier, Siming Bayer, and Prajol Shrestha:
Pattern Recognition Lab, Department of Computer Science,
Friedrich-Alexander-Universit\"at Erlangen-N\"urnberg,
Martensstr.\ 3, 91058 Erlangen, Germany.

\noindent Isa Mammadli and Siddhant Mohapatra:
PULS Group, Institute for Theoretical Physics,
Friedrich-Alexander-Universit\"at Erlangen-N\"urnberg,
Cauerstr.\ 3, 91058 Erlangen, Germany.

\noindent M. Delens, M. M\'elard, Maxime Noirhomme,
N. Vanesse, F. W\'ery, and Nicolas Vandewalle:
GRASP, Institute of Physics B5a,
University of Li\`ege, B-4000 Li\`ege, Belgium.

\noindent Filip Novkoski:
PULS Group, Institute for Theoretical Physics,
Friedrich-Alexander-Universit\"at Erlangen-N\"urnberg,
Cauerstr.\ 3, 91058 Erlangen, Germany;
and GRASP, Institute of Physics B5a,
University of Li\`ege, B-4000 Li\`ege, Belgium.

\noindent Jayant Pande:
Department of Physical and Natural Sciences,
FLAME University, Pune, India.

\noindent Ana-Sun\v{c}ana Smith:
PULS Group, Institute for Theoretical Physics,
Friedrich-Alexander-Universit\"at Erlangen-N\"urnberg,
Cauerstr.\ 3, 91058 Erlangen, Germany;
and Group for Computational Life Sciences,
Division of Physical Chemistry,
Rudjer Bo\v{s}kovi\'c Institute,
10000 Zagreb, Croatia.

\noindent\textbf{Corresponding author:}
[Name and email address]

\section*{Abstract}

Robotic active matter provides a new experimental route to
nonequilibrium statistical physics. Unlike many biological
or chemically driven active systems, robotic collectives
permit researchers to modify propulsion, interaction
rules, sensing, communication, feedback, and mechanical
coupling deliberately. Experiments range from macroscopic
robot swarms and vibration-driven particles to programmable
microscopic agents, feedback-controlled colloids, and
active mechanical metamaterials. These systems exhibit
collective ordering, clustering, transport, synchronization,
self-assembly, and nonreciprocal mechanical response.

We propose a critical, PRISMA 2020-compliant living
systematic review of this emerging experimental field.
The October 17, 2026 evidence snapshot comprises
489 included studies. We organize the literature around
five questions: Which robotic platforms realize active
matter? Which collective phenomena have been observed?
Which physical theories have been tested quantitatively?
How do sensing and control modify collective behavior?
And how reproducible is the experimental evidence?

Our central distinction is between demonstrating a
collective pattern and independently validating its
physical explanation. We assess theory--experiment
agreement, finite-size effects, replication, measurement
validity, and data and code availability. The review
integrates physical interpretation with an auditable
evidence synthesis rather than presenting a catalogue
of robotic technologies.

A companion public website provides a searchable,
reconfigurable taxonomy of the included studies.
Monthly updates and reviewed community submissions
through GitHub maintain the evidence base as the
field develops.

\section*{1. Motivation and timeliness}

Active-matter physics seeks to understand how
energy-consuming constituents generate collective
behavior far from equilibrium. Its theories connect
microscopic propulsion and interactions to macroscopic
ordering, transport, fluctuations, and mechanical
response. Yet, the relationship between idealized
models and experimental systems remains difficult
to establish.

Robotic active matter offers a distinctive solution
to this experimental challenge. Interactions can
often be programmed or modified directly, allowing
the assumptions of a model to become independently
adjustable experimental variables.

Macroscopic robots have demonstrated collective
alignment and distributed coordination
\cite{Turgut2008selforganizedflocka,Vsrhelyi2018optimizedflockingo,Zhou2024anticipationpromote}.
Vibration-driven robots have enabled experiments
on persistent motion, aggregation, and
interaction-mediated organization
\cite{Welch2018aggregatesuppressio,Deblais2019howtoforceanarmy,Chvykov2020lowrattlingapred}.
Microscopic robotic assemblies and
feedback-controlled colloids provide access
to field-mediated and programmable interactions
\cite{Xie2020bioinspiredsoftmic,Pauer2021programmabledesign,Buerle2020formationofstable,Heuthe2025tunablecolloidalsw}.
Active robotic structures extend these
questions to mechanical response and
nonreciprocal elasticity
\cite{Brandenbourger2021limitcyclesturnac,Coulais2025activemetamaterials}.

These advances have created an identifiable
experimental research field at the intersection
of active-matter physics, robotics, and
programmable materials. However, its results
remain dispersed across physics, engineering,
and materials-science publications.

The timing is therefore appropriate for
a comprehensive review that establishes
what is experimentally known and separates
robust physical conclusions from
platform-specific demonstrations.

The literature has also reached sufficient
breadth to support systematic evidence mapping.
Our corrected evidence snapshot contains
489 studies, spanning multiple physical scales,
interaction mechanisms, and theoretical
frameworks.

Yet, numerical breadth alone does not establish
scientific maturity. The evidence includes
many qualitative comparisons and fewer
low-risk quantitative tests. Only 81 studies
receive low-risk ratings in the review's
quantitative theory-comparison domain,
and 38 receive low-risk ratings for
data and code availability.

These findings indicate a timely opportunity
for critical assessment. The field needs
a synthesis that identifies experimentally
supported principles and the measurements
required to test them more decisively.

\section*{2. Relationship to existing reviews}

The proposed article bridges three established
review traditions.

First, general active-matter reviews develop
the theoretical language of self-propulsion,
symmetry breaking, hydrodynamic fields,
fluctuations, and nonequilibrium phases
\todocite{Ramaswamy, 2010, mechanics and
statistics of active matter}
\todocite{Marchetti et al., 2013, hydrodynamics
of soft active matter}
\todocite{Bechinger et al., 2016, active
particles in complex and crowded environments}.

Second, swarm-robotics reviews emphasize
distributed algorithms, sensing, communication,
collective tasks, and engineering scalability
\todocite{Brambilla et al., 2013, swarm robotics review}
\todocite{review of swarm-robotics methods
and distributed collective control}.

Third, platform-specific reviews examine
microrobotic collectives, active colloids,
and programmable mechanical materials
\todocite{review of collective magnetic microrobotics}
\todocite{review of active mechanical metamaterials}.

Our proposed article occupies a different
position. It takes physical robotic experiments
as the common evidence base and asks what
they establish about active-matter theories.

A robotic swarm that successfully navigates
as a formation may be an important engineering
achievement, but its task performance alone
does not establish a collective physical law.
Conversely, an experiment that measures
orientation correlations, fluctuations,
phase-dependent transport, or mechanical
response can provide direct constraints
on a theoretical model.

This distinction is the basis of our
eligibility criteria and critical synthesis.

\section*{3. Scope and review objectives}

The review includes primary studies of at least
two interacting, self-propelled or actively
driven artificial agents in physical hardware,
with at least one quantitative collective
or interaction observable.

The scope encompasses conventional mobile robots,
vibration-driven systems, smarticles, modular
assemblies, robotic micro- and nano-agents,
feedback-controlled active colloids, and
robotic active solids.

Purely computational studies without the
specified experimental hardware connection
are excluded. Engineering studies reporting
only task-performance measures are also
excluded when they provide no qualifying
collective physical observable.

We address five prespecified objectives:

\noindent\textbf{O1---Platforms.}
Map robotic implementations of active matter
and identify the physical interactions and
observables accessible to each platform.

\noindent\textbf{O2---Phenomena.}
Identify collective behaviors observed
experimentally, distinguishing their
phenomenological similarities from
their underlying mechanisms.

\noindent\textbf{O3---Theory.}
Determine which active-matter frameworks
have been compared quantitatively with
robotic experiments and assess
the reported agreement.

\noindent\textbf{O4---Control.}
Characterize feedback, sensing,
communication, and learning as
experimentally adjustable ingredients
of collective dynamics.

\noindent\textbf{O5---Reproducibility.}
Assess experimental replication,
finite-size and boundary effects,
measurement validity,
agent heterogeneity,
and data and code availability.

The review follows PRISMA 2020
\cite{Page2021prisma} and applies
a documented search, screening,
extraction, adjudication, and
risk-of-bias workflow.

\section*{4. Proposed table of contents}

\begin{enumerate}
\item Introduction
\item Methods
\item Overview of the evidence base
\item Robotic platforms as active matter
\item Collective phenomena in robotic active matter
\item Confronting theory with robotic experiments
\item Control, sensing, and programmability
\item Discussion
\item Outlook
\item Conclusions
\end{enumerate}

\section*{5. Organization and scientific argument}

\subsection*{Introduction and Methods}

The Introduction motivates robotic collectives
as physical realizations of active-matter models,
distinguishes collective demonstrations from
theory tests, and establishes the five objectives.

The Methods describes the systematic search,
eligibility decisions, extraction taxonomy,
LLM-assisted screening and adjudication,
validation procedures, and six-domain
risk-of-bias assessment.

The methodological emphasis is important:
experimental diversity does not justify
treating every study as equally informative,
and a reported quantitative comparison
does not necessarily establish
independent prediction.

\subsection*{Evidence base and experimental platforms}

The evidence overview documents study selection,
publication trends, hardware characteristics,
and the distribution of experimental approaches.

The platforms section then examines each
experimental class in physical terms:
what generates activity, how agents interact,
which quantities can be measured, and
what limits theoretical inference.

The discussion connects traditional robot
swarms to mechanically driven systems,
microscopic assemblies, feedback-controlled
particles, and active structures
\cite{Rubenstein2014programmableselfas,Li2019particleroboticsba,Savoie2019arobotmadeofrobo,Sun2021swarmingmicrodrople,Yang2022emergentmicroroboti}.

This platform-based analysis establishes
the experimental capabilities on which
the following physical synthesis depends.

\subsection*{Collective phenomena}

The phenomena section reorganizes the same
evidence according to collective observables:
flocking, clustering, swarming and vortices,
jamming, chirality, active elasticity,
transport, synchronization, self-assembly,
and boundary accumulation.

Its central question is whether apparently
similar collective states are supported
by comparable physical mechanisms.

For example, clustering may be produced
through programmed attraction, steric
interactions, or field-mediated assembly.
The section distinguishes these possibilities
rather than treating all aggregates as
equivalent evidence for phase separation.

Likewise, active rotation can arise from
externally imposed chirality or
interaction-dependent collective response
\cite{Fruchart2022nonreciprocalphase,Chen2024persistentandrespo,Casiulis2025ageometricconditio}.

This provides the bridge from observed
phenomena to explanatory theory.

\subsection*{Theory--experiment confrontation}

The theory section constitutes the
principal physics-focused contribution.

It compares Vicsek-type descriptions,
active Brownian particle models,
run-and-tumble dynamics,
kinetic and hydrodynamic theories,
active elasticity,
other statistical-mechanical approaches,
and agent-based simulations.

The analysis distinguishes three
levels of evidence:
physical implementation of a rule,
agreement with measured observables,
and prediction under independently
constrained parameters.

This distinction is particularly important
for agent-based simulation, which frequently
reproduces robotic controllers directly.

The section identifies experiments
that provide quantitative physical
constraints, including measurements
of active transport, collective response,
and mechanical behavior
\cite{Tarcai2011patternstransition,Barone2024experimentalandnum,Zion2023morphologicalcomput,Binysh2026moreislessinunpe}.

It also explains why numerical agreement
obtained through extensive fitting
is less decisive than successful
prediction of independent observables.

\subsection*{Control and programmability}

The control section examines passive rules,
on-board control, external feedback,
centralized control, and learned policies.

Sensing and communication are treated
as physical interaction mechanisms,
not merely as engineering infrastructure.

The analysis asks which effective
interactions are created by
information processing and which
collective effects remain predictions.

This perspective connects programmable
robotics directly to questions of
interaction symmetry, nonreciprocity,
delay, noise, and collective stability.

\subsection*{Discussion, Outlook, and Conclusions}

The Discussion integrates the findings
across all five objectives and assigns
claim-specific certainty levels.

It examines limitations in independent
replication, finite-size analysis,
data availability, and quantitative
theory tests, while also evaluating
the review's own methodological limitations.

The Outlook proposes decisive experiments:
parameter-constrained alignment tests,
controlled comparisons of aggregation
mechanisms, nonreciprocity measurements,
and replicated mechanical-response studies.

The Conclusions summarize what robotic
active matter has established and
which broader theoretical claims
remain insufficiently tested.

\section*{6. Living-review infrastructure}

The article will be maintained as a living
systematic review, following the established
methodology for continually updated
evidence syntheses
\cite{Elliott2014living,Elliott2017living}.

Database searches will be repeated monthly.
New studies will undergo the documented
screening and extraction procedure,
and changes will be recorded in
dated evidence snapshots.

A companion public website provides
a searchable, reconfigurable taxonomy
following the interactive living-review
concept of Matl et al.~\cite{Matl2017vascular}.

Readers can choose which extracted
attributes form the successive levels
of the taxonomy tree and reorder them.
Thus, the same literature can be explored
by platform, phenomenon, interaction
modality, control architecture,
theoretical framework, or
theory--experiment agreement.

This extends the online-updated survey
model used in other research areas
\cite{Litjens2017survey}.
The structured review of brain-shift compensation
by Bayer et al.~\cite{Bayer2017brainshift} provides
another example of a technique-oriented review
whose taxonomy organizes a heterogeneous
methodological literature.

Community participation is supported
through reviewed GitHub pull requests.
Anyone may propose a missing paper
by adding one small YAML file.
Automated checks verify the DOI
against Crossref and flag records
already held in the review.

Maintainers review and merge candidate
entries. The next monthly update
screens them with the same panel
and eligibility criteria used
for database records.

Such contributions are documented
as records identified through
other methods under PRISMA 2020
\cite{Page2021prisma}.

This design provides a practical
mechanism for keeping the review
current without relaxing its
evidence-selection standards.

\section*{7. Selected key references}

The following references illustrate the
range of experimental evidence and
review methodology. They are organized
by scientific role rather than
publication date.

\noindent\textbf{Collective motion and robotic swarms:}
\cite{Turgut2008selforganizedflocka,Tarcai2011patternstransition,Rubenstein2014programmableselfas,Vsrhelyi2018optimizedflockingo,Zhou2024anticipationpromote}.

\noindent\textbf{Mechanically driven active systems:}
\cite{Welch2018aggregatesuppressio,Deblais2019howtoforceanarmy,Chvykov2020lowrattlingapred,Barone2024experimentalandnum,Wang2025hyperuniformnetwork}.

\noindent\textbf{Microscopic and feedback-controlled collectives:}
\cite{Xie2020bioinspiredsoftmic,Pauer2021programmabledesign,Sun2021swarmingmicrodrople,Buerle2020formationofstable,Heuthe2025tunablecolloidalsw}.

\noindent\textbf{Modular robots and active mechanical response:}
\cite{Li2019particleroboticsba,Savoie2019arobotmadeofrobo,Brandenbourger2021limitcyclesturnac,Coulais2025activemetamaterials,Binysh2026moreislessinunpe}.

\noindent\textbf{Systematic and living-review methodology:}
\cite{Page2021prisma,McGowan2016press,Elliott2014living,Elliott2017living,Matl2017vascular,Litjens2017survey,Bayer2017brainshift}.

\section*{8. Author expertise and suitability}

The proposed author team combines expertise in
statistical physics, active matter, robotic
experiments, computational methods, and
machine learning.

Andreas Maier heads the Pattern Recognition
Lab at FAU and works on medical image
processing, machine learning, and pattern
recognition for health engineering. He is
a co-author of the GRASPion and brainbot
digital-twin papers.

Siming Bayer is a research-group head
at the Pattern Recognition Lab, FAU,
with a background in medical image fusion.

M. Delens and M. M\'elard contributed
to the GRASPion open-source bristlebot
paper at GRASP, University of Li\`ege.

Isa Mammadli is a student assistant
in the PULS group and a co-developer
of the brainbot platform. Mammadli
is first author of the brainbot
digital-twin and on-board-control
paper.

Siddhant Mohapatra is a research associate
in the PULS group at FAU and is listed
in the DFG Research Training Group
``Fracture across scales.''

Maxime Noirhomme is first author
of the original brainbot paper
and a co-author of the GRASPion
paper.

Filip Novkoski works on programmable
bristlebots and brainbots for active
matter, has affiliations with PULS
and GRASP, and is first author
of the GRASPion paper.

Jayant Pande is an Assistant Professor
of Physics at FLAME University whose
research interests include complex-systems
physics and theoretical ecology.
He is a corresponding author
of the brainbot papers.

Prajol Shrestha is a doctoral researcher
at the Pattern Recognition Lab working
on multimodal signal processing with AI
and predictive control and flocking
of vibration-driven robots.

Ana-Sun\v{c}ana Smith is Professor
of Theoretical Physics at FAU and head
of the PULS group, with expertise
in biological membranes, cell adhesion,
and soft and active matter.

N. Vanesse is a co-author of the
original brainbot paper. F. W\'ery
is a co-author of the GRASPion paper.

Nicolas Vandewalle is Professor of Physics
and director of GRASP at the University
of Li\`ege, with expertise in statistical
physics of soft and active matter.
He is the last author of the
three brainbot papers.

The interdisciplinary nature of the
article requires expertise spanning
physical theory, experimental hardware,
and computational methods. These
complementary backgrounds support
the proposed critical synthesis.

\section*{9. Contribution to Reviews of Modern Physics}

We propose this article for
\textit{Reviews of Modern Physics}
because robotic active matter now
poses substantive questions about
nonequilibrium physics that cannot
be addressed through a platform
catalogue or an engineering survey.

The review will establish which
collective behaviors are experimentally
reproducible, which models have
quantitative support, and where
theoretical explanations remain
limited by fitting, finite size,
or incomplete characterization.

Its distinguishing contribution
is the connection between
physical mechanisms and
evidence quality.

We will not infer universal
behavior from similar-looking
collective patterns or treat
successful controller execution
as independent theory validation.

Instead, the review will
identify the observables
that constrain models,
the experiments that
distinguish competing
mechanisms, and the
conditions under which
a theoretical description
can be generalized.

The systematic methodology
makes these judgments
auditable. The living
infrastructure allows
them to be revised
as stronger evidence
becomes available.

For the physics community,
the result will be both
a critical account of
an emerging experimental
field and a guide to
the decisive tests
needed for its future
development.

The central message
is that robotic
programmability becomes
most scientifically
powerful when it
makes theoretical
assumptions experimentally
testable, rather than
merely making collective
behavior possible.
